\section{Clean \agnews{}}
\label{sec:results-agnews-clean}

\input{tables/nlp_agnews_clean}

The clean \agnews{} experiment provides a stable Transformer fine-tuning setting where \adamw{} with warmup-linear scheduling is expected to be a strong baseline. As shown in \Cref{tab:agnews-clean}, this expectation is confirmed. Flat \adamw{} reaches $93.98\%$ peak validation accuracy and $93.52\%$ final accuracy, while \adamw{} with warmup-linear scheduling improves this to $94.44\%$ peak accuracy and $94.36\%$ final accuracy.

\method{} improves over flat \adamw{} in peak validation accuracy, reaching $94.40\%$, but its final accuracy of $94.12\%$ remains below the warmup-linear baseline. Combining \method{} with warmup-linear scheduling gives $94.38\%$ peak accuracy and $94.34\%$ final accuracy, which is essentially matched with \adamw{}+warmup-linear. The best-to-final drop is small for both scheduled configurations, with $0.09$ percentage points for \adamw{}+warmup-linear and $0.04$ percentage points for \method{}+warmup-linear.

These results show that \method{} does not provide a clear improvement over a strong standard scheduler in this clean fine-tuning setting. Instead, clean \agnews{} acts as a boundary case: adaptive control can match the warmup-linear baseline, but most of the available gain over flat \adamw{} is already captured by the conventional scheduler. This supports the broader interpretation that \method{} is most useful when the training regime leaves room for adaptive exploration, rather than in stable settings where standard fine-tuning recipes are already well matched to the task.